El conjunto de casos de prueba utilizado para esta tarea fue generado con el objetivo de analizar el comportamiento práctico de los algoritmos en entradas variadas y controladas. Se definieron $K = 10$ casos, con pares de cadenas de texto $(s, t)$ que varían en tamaño total desde 10 hasta 320 caracteres. Los strings fueron diseñados para asegurar una mezcla de patrones repetitivos, diferencias parciales y presencia de espacios, lo cual permite evaluar distintos tipos de dificultad para los algoritmos.

Los tamaños fueron elegidos de forma incremental, comenzando con secuencias cortas ($|s| = |t| = 5$) y aumentando progresivamente hasta cadenas de tamaño 160. Para cada par, se mantuvo una cierta similitud estructural entre $s$ y $t$, permitiendo que el algoritmo encuentre subsecuencias comunes sin ser triviales. Algunos ejemplos incluyen cadenas con diferencias pequeñas, diferencias múltiples, cadenas muy similares y otras más divergentes.

\begin{itemize}
    \item Casos 1-2: cadenas cortas con diferencias simples.
    \item Casos 3-4: cadenas con espacios y bloques modificados parcialmente.
    \item Casos 5-10: cadenas largas con patrones estructurados, diseñadas para poner a prueba la escalabilidad y eficiencia de los algoritmos.
\end{itemize}

Los casos fueron generados mediante un script en Python que garantiza la reproducibilidad de los experimentos. Este script fue ejecutado una sola vez para producir un archivo \texttt{a.txt}, que luego se utilizó como entrada para ambos algoritmos. Así se asegura una comparación justa.

Debido a la naturaleza exponencial del enfoque por fuerza bruta, este solo fue ejecutado hasta el segundo caso del conjunto. A partir del tercer caso ($|s| = |t| = 20$), el tiempo de ejecución superó las dos horas sin finalizar, por lo que se omitieron esos resultados. En cambio, el enfoque de programación dinámica pudo ejecutarse completamente en todos los casos, demostrando su superioridad en términos de escalabilidad.

Este dataset, al ser estático y variado, permite un análisis claro del rendimiento de ambos enfoques y refleja adecuadamente la motivación de la tarea: contrastar paradigmas algorítmicos bajo condiciones equivalentes.
